% W4i32_QG_Core.tex
% DFD 17-Agent Research Programme — Iteration 32 (i32)
% Agents 1 (Formalism/QG Lead), 4 (Perturbation/Loop Corrections),
%         7 (Superposition Problem), 9 (Renormalization)
% FIRST ITERATION of new 20-cycle (i32-i51): SOLVE DFD's QUANTUM GRAVITY
% Date: 2026-03-22

\documentclass[11pt,a4paper]{article}
\usepackage[utf8]{inputenc}
\usepackage{amsmath,amssymb,amsfonts,amsthm}
\usepackage{hyperref}
\usepackage{booktabs}
\usepackage{longtable}
\usepackage{geometry}
\usepackage{xcolor}
\usepackage{slashed}
\usepackage{bbm}
\geometry{margin=2.5cm}

\newtheorem{theorem}{Theorem}
\newtheorem{proposition}{Proposition}
\newtheorem{lemma}{Lemma}
\newtheorem{corollary}{Corollary}
\newtheorem{definition}{Definition}
\newtheorem{remark}{Remark}

\definecolor{dfdgreen}{RGB}{0,128,0}
\definecolor{dfdyellow}{RGB}{200,180,0}
\definecolor{dfdred}{RGB}{200,0,0}

\newcommand{\GREEN}{\textcolor{dfdgreen}{\textbf{GREEN}}}
\newcommand{\YELLOW}{\textcolor{dfdyellow}{\textbf{YELLOW}}}
\newcommand{\RED}{\textcolor{dfdred}{\textbf{RED}}}

\title{DFD i32: Quantum Gravity Core\\[6pt]
\large Agents 1 (Formalism/QG Lead), 4 (Loop Corrections),\\
7 (Superposition Problem), 9 (Renormalization)\\[6pt]
\normalsize Iteration 1/20 of Quantum Gravity Cycle (i32--i51)}
\author{DFD 17-Agent Research Programme}
\date{2026-03-22}

\begin{document}
\maketitle
\tableofcontents
\newpage

%=============================================================================
\section{Agent 1: Formalism / Quantum Gravity Lead}
\label{sec:agent1}
%=============================================================================

\subsection{Mandate}

The central question: Can the DFD scalar field $\psi$ be quantized as a scalar QFT on flat Minkowski space? Write the path integral. Determine the propagator. Assess renormalizability. Identify vertices from the nonlinear k-essence Lagrangian.

\subsection{New Literature (i32)}

\begin{enumerate}
\item \textbf{Takahashi \& Motohashi (2026).} arXiv: \href{https://arxiv.org/abs/2601.09164}{2601.09164}. Third-order disformal Horndeski --- establishes the complete classification of disformal scalar-tensor theories. DFD's k-essence sector is a proper subclass. \textbf{Grade: A.}

\item \textbf{UV completions of scalar-tensor EFTs.} Blas, Criado \& Noumi (2025). arXiv: \href{https://arxiv.org/abs/2507.11426}{2507.11426}. JHEP (Jan 2026). One-loop matching of scalar-tensor EFTs to UV completions with massive spin-0, 1/2, 1 particles. Demonstrates that shift-symmetric scalar Gauss-Bonnet always breaks shift symmetry at one loop. \textbf{Directly relevant to DFD's shift symmetry.} \textbf{Grade: A.}

\item \textbf{Gravitationally induced UV completion of O(N) scalar theory.} arXiv: \href{https://arxiv.org/abs/2601.20820}{2601.20820} (Jan 2026). Wilsonian FRG in proper-time formulation for scalar field non-minimally coupled to gravity. Demonstrates asymptotic safety of scalar potential + non-minimal coupling. \textbf{Grade: B.}

\item \textbf{Nonlocal corrections to scalar field effective action in de Sitter.} arXiv: \href{https://arxiv.org/abs/2601.22644}{2601.22644} (Jan 2026). Schwinger-Keldysh one-loop effective action for scalar field in FLRW. Memory term + stochastic noise. \textbf{Grade: B.}

\item \textbf{Renormalising the field-space geometry.} JHEP (Jul 2025). One-loop quantum corrections to scalar EFTs from geometric viewpoint. Field-space curvature receives momentum- and scale-dependent renormalization. \textbf{Grade: A.}

\item \textbf{Gradient corrections to the quantum effective action.} arXiv: \href{https://arxiv.org/abs/2208.12142}{2208.12142v3} (Sep 2024). Two-loop gradient expansion for scalar effective action. Single-field gradient corrections vanish at one loop but appear at two loops. \textbf{Grade: B.}
\end{enumerate}

\subsection{The DFD Classical Action}

We begin with the DFD scalar-sector action in relativistic (Lorentzian) form:
\begin{equation}\label{eq:DFD-action}
S_\psi = \int d^4x \left\{
  \frac{a_0^2}{8\pi G c^4}\,\mathcal{L}_{\rm kin}(\chi)
  + \frac{1}{2}\psi\, T
\right\},
\end{equation}
where $T = T^\mu{}_\mu / c^2$ is the trace of the matter stress-energy tensor (divided by $c^2$ for dimensional convenience), and the kinetic function is
\begin{equation}\label{eq:L-kin}
\mathcal{L}_{\rm kin}(\chi) = \chi - \ln(1+\chi), \qquad
\chi \equiv \frac{(\partial\psi)^2}{a_\star^2},
\end{equation}
with $(\partial\psi)^2 = \eta^{\mu\nu}\partial_\mu\psi\,\partial_\nu\psi$ (signature $-,+,+,+$) and $a_\star = 2a_0/c^2$.

This is a $P(X)$ theory with $P(X) = a_0^2/(8\pi G c^4) \cdot [\chi - \ln(1+\chi)]$ where $X = (\partial\psi)^2$ and $\chi = X/a_\star^2$.

\subsection{Expansion of the Kinetic Lagrangian}

Expanding $\mathcal{L}_{\rm kin}(\chi)$ around $\chi = 0$:
\begin{equation}\label{eq:L-expansion}
\mathcal{L}_{\rm kin}(\chi) = \chi - \ln(1+\chi)
= \sum_{n=2}^{\infty} \frac{(-1)^n}{n}\,\chi^n
= \frac{\chi^2}{2} - \frac{\chi^3}{3} + \frac{\chi^4}{4} - \cdots
\end{equation}

\textbf{Critical observation}: $\mathcal{L}_{\rm kin}$ starts at $\chi^2$, not $\chi$. This is consistent with the established result $W'(0) = 0$ (confirmed i20). The field $\psi$ has \textbf{no standard kinetic term} at the quadratic level.

Substituting $\chi = X/a_\star^2$:
\begin{equation}\label{eq:L-X-expansion}
\mathcal{L}_{\rm kin} = \frac{X^2}{2a_\star^4} - \frac{X^3}{3a_\star^6} + \frac{X^4}{4a_\star^8} - \cdots
\end{equation}

In terms of the canonically dimensioned field ($[\psi]$ is dimensionless in DFD):
\begin{equation}
\mathcal{L}_{\rm kin} = \frac{(\partial\psi)^4}{2a_\star^4}
- \frac{(\partial\psi)^6}{3a_\star^6} + \cdots
\end{equation}

\subsection{The Propagator Problem}
\label{sec:propagator}

\begin{theorem}[DFD Propagator Degeneracy]
The DFD scalar field $\psi$ has no standard free-field propagator. The inverse propagator vanishes identically in the free-field limit:
\begin{equation}
\mathcal{L}_{\rm free} = \lim_{\psi\to 0} \mathcal{L}_{\rm kin} = 0.
\end{equation}
There is no $(\partial\psi)^2$ term in the Lagrangian; the leading kinetic term is $(\partial\psi)^4$.
\end{theorem}

\begin{proof}
From Eq.~(\ref{eq:L-expansion}), $\mathcal{L}_{\rm kin}(\chi) = \chi^2/2 + O(\chi^3)$. Since $\chi = (\partial\psi)^2/a_\star^2$, the leading term in the action is $\sim (\partial\psi)^4$. The standard propagator is defined from the quadratic part of the action $\sim (\partial\psi)^2$, which is \emph{absent}. Therefore, the free-field propagator $\langle \psi(x)\psi(y)\rangle_{\rm free}$ does not exist as a standard Feynman propagator.
\end{proof}

\begin{remark}
This is \textbf{not} a deficiency --- it is a deep structural feature. DFD's $\psi$ is not a propagating quantum field in the free-field sense. It is a \textbf{constraint field} analogous to the Coulomb potential in electrostatics (which also has no $\dot{A}_0^2$ kinetic term in the Lagrangian). The field $\psi$ mediates interactions but does not have free-particle excitations at the linearized level.
\end{remark}

\subsection{Quantization Strategy: Background Field Method}
\label{sec:background-field}

Since $\psi$ has no free propagator, we must quantize around a \emph{non-trivial background}. Let:
\begin{equation}\label{eq:background-split}
\psi(x) = \bar{\psi}(x) + \varphi(x),
\end{equation}
where $\bar{\psi}$ solves the classical field equation (sourced by matter) and $\varphi$ is the quantum fluctuation.

Expanding $\mathcal{L}_{\rm kin}$ to second order in $\varphi$ around $\bar{\psi}$:

Define $\bar{\chi} = (\partial\bar{\psi})^2/a_\star^2$ and $\bar{X} = (\partial\bar{\psi})^2$. The second-order expansion gives:
\begin{equation}\label{eq:second-order}
\mathcal{L}_{\rm kin}^{(2)} = P_{,X}\,(\partial\varphi)^2
+ 2P_{,XX}\,(\partial_\mu\bar{\psi}\,\partial^\mu\varphi)^2,
\end{equation}
where $P_{,X} = dP/dX$ and $P_{,XX} = d^2P/d X^2$, evaluated at $\bar{X}$.

Computing the derivatives of $P(\chi) = \chi - \ln(1+\chi)$:
\begin{align}
P_{,\chi} &= 1 - \frac{1}{1+\chi} = \frac{\chi}{1+\chi} = \mu(\chi), \label{eq:P-chi}\\
P_{,\chi\chi} &= \frac{1}{(1+\chi)^2}. \label{eq:P-chichi}
\end{align}

Converting to $X$ derivatives ($\chi = X/a_\star^2$, $P_X = P_\chi / a_\star^2$, $P_{XX} = P_{\chi\chi}/a_\star^4$):
\begin{align}
P_{,X} &= \frac{1}{a_\star^2}\,\mu(\bar{\chi}), \\
P_{,XX} &= \frac{1}{a_\star^4}\,\frac{1}{(1+\bar{\chi})^2}.
\end{align}

The quadratic Lagrangian for fluctuations becomes:
\begin{equation}\label{eq:L2-final}
\boxed{
\mathcal{L}^{(2)} = \frac{1}{a_\star^2}\left[
\mu(\bar\chi)\,(\partial\varphi)^2
+ \frac{2\bar{X}}{a_\star^2(1+\bar\chi)^2}\,
(\hat{n}^\mu\partial_\mu\varphi)^2
\right],
}
\end{equation}
where $\hat{n}^\mu = \partial^\mu\bar{\psi}/|\partial\bar{\psi}|$ is the unit gradient direction of the background.

\begin{proposition}[Background-Dependent Propagator]
\label{prop:propagator}
In a non-trivial background $\bar{\psi}$ with gradient $|\nabla\bar{\psi}| \neq 0$, the $\psi$-fluctuation $\varphi$ acquires an effective propagator. In momentum space, for a homogeneous background with $\partial_\mu\bar\psi = (\dot{\bar\psi}, \nabla\bar\psi)$:
\begin{equation}\label{eq:propagator}
G_\varphi^{-1}(k) = \frac{2}{a_\star^2}\left[
\mu(\bar\chi)\,k^2
+ \frac{2\bar{X}}{a_\star^2(1+\bar\chi)^2}\,(\hat{n}\cdot k)^2
\right].
\end{equation}
The propagator is \textbf{anisotropic} --- it depends on the direction of $k$ relative to the background gradient $\nabla\bar\psi$.
\end{proposition}

\subsection{Sound Speed of $\psi$-Fluctuations}
\label{sec:sound-speed}

The sound speed for $\varphi$-fluctuations propagating in the background $\bar\psi$ is:
\begin{equation}\label{eq:cs-general}
c_s^2 = \frac{P_{,X}}{P_{,X} + 2X\,P_{,XX}}
= \frac{\mu(\bar\chi)}{\mu(\bar\chi) + 2\bar\chi/(1+\bar\chi)^2}.
\end{equation}

Simplifying using $\mu(\chi) = \chi/(1+\chi)$:
\begin{equation}
c_s^2 = \frac{\chi/(1+\chi)}{\chi/(1+\chi) + 2\chi/(1+\chi)^2}
= \frac{1}{1 + 2/(1+\chi)} = \frac{1+\chi}{3+\chi}.
\end{equation}

\begin{theorem}[Sound Speed Bounds]
\label{thm:sound-speed}
For the DFD k-essence theory with $\mathcal{L}_{\rm kin} = \chi - \ln(1+\chi)$:
\begin{equation}
c_s^2(\chi) = \frac{1+\chi}{3+\chi},
\end{equation}
with the limits:
\begin{align}
\chi \to 0 &\quad(\text{deep MOND}): & c_s^2 &\to \frac{1}{3}, \\
\chi \to \infty &\quad(\text{Newtonian}): & c_s^2 &\to 1.
\end{align}
The sound speed is everywhere \textbf{subluminal}: $1/3 \leq c_s^2 \leq 1$.
\end{theorem}

\begin{remark}
The $c_s^2 \to 1/3$ in the deep-MOND limit is significant: it means $\psi$-quantum fluctuations propagate at $c/\sqrt{3}$ in weak-field regions. This is the conformal value, suggesting a hidden conformal symmetry in the deep-MOND sector.
\end{remark}

\textbf{CRITICAL RECONCILIATION}: The $c_s^2 \to 0$ theorem from earlier iterations (W4i15) was derived under the \emph{linearized FRW background} where $\bar\chi \to 0$. The present result $c_s^2 \to 1/3$ applies to the fully nonlinear AQUAL equation on a non-trivial background. The resolution: on a truly homogeneous FRW background $\nabla\bar\psi = 0$, so the entire kinetic sector is degenerate. On any physical (inhomogeneous) background, $\bar\chi > 0$ and $c_s^2 \geq 1/3$. The ``$c_s = 0$'' result was an artifact of the degenerate FRW limit, as recognized at i20.

\subsection{Path Integral Formulation}
\label{sec:path-integral}

The DFD partition function in the background field formalism:
\begin{equation}\label{eq:path-integral}
Z[\bar\psi, J] = \int \mathcal{D}\varphi\,\exp\left\{
\frac{i}{\hbar}\int d^4x\left[
\frac{a_0^2}{8\pi G c^4}\,\mathcal{L}_{\rm kin}[\bar\psi+\varphi]
+ \frac{1}{2}(\bar\psi+\varphi)T + J\varphi
\right]\right\}.
\end{equation}

The overall coupling constant is:
\begin{equation}
\frac{a_0^2}{8\pi G c^4} = \frac{a_0^2}{8\pi G c^4}
\approx \frac{(1.2\times 10^{-10})^2}{8\pi \times 6.67\times 10^{-11} \times (3\times 10^8)^4}
\approx 1.07 \times 10^{-42}\;\text{m}^{-2}\text{s}^2.
\end{equation}

This extraordinarily small prefactor means that the action is dominated by the matter coupling $\psi T/2$, and quantum fluctuations of $\psi$ are \textbf{extremely suppressed} relative to classical solutions.

\subsection{Effective Planck Scale for $\psi$}
\label{sec:psi-planck}

Define the effective coupling for $\psi$-matter interaction. The $\psi$-field couples to matter through:
\begin{equation}
\mathcal{L}_{\rm int} = \frac{1}{2}\psi\,T.
\end{equation}

The dimensionless coupling strength, compared to the kinetic sector normalization, defines an effective Planck mass for $\psi$:
\begin{equation}\label{eq:M-psi-Planck}
M_{\psi,\rm Pl}^2 \equiv \frac{a_0^2}{4\pi G c^4}
= \frac{2a_\star^2 c^4}{4\pi G \cdot 4}
= \frac{a_\star^2 c^4}{8\pi G}.
\end{equation}

Numerically:
\begin{equation}
M_{\psi,\rm Pl} \approx 1.16 \times 10^{-21}\;\text{kg}
\approx 6.5 \times 10^{5}\;\text{eV}/c^2.
\end{equation}

This is \textbf{seventeen orders of magnitude below the Planck mass} $M_{\rm Pl} = 2.18 \times 10^{-8}$ kg. The $\psi$-sector ``goes quantum'' at $\sim 0.65$ MeV --- coincidentally near the electron mass scale.

\begin{theorem}[DFD Quantum Gravity Scale]
\label{thm:QG-scale}
The scale at which quantum effects of the $\psi$-field become important is:
\begin{equation}
\Lambda_\psi = \sqrt{\frac{a_0^2}{G c^2}} = a_0\,\sqrt{\frac{1}{G c^2}} \approx 0.65\;\text{MeV}.
\end{equation}
This is set by $a_0$ (the MOND acceleration) and $G$ (Newton's constant), and is independent of $\hbar$. Quantum $\psi$-effects are relevant at \textbf{nuclear and atomic scales}, not at the Planck scale.
\end{theorem}

\begin{remark}
This is a \textbf{radical departure} from standard quantum gravity. In GR, quantum gravity effects are suppressed by $E/M_{\rm Pl} \sim E/(10^{19}\;\text{GeV})$. In DFD, quantum $\psi$-effects are suppressed by $E/\Lambda_\psi \sim E/(0.65\;\text{MeV})$. However, this does NOT mean DFD predicts large quantum gravity effects at MeV scales, because the $\psi$-fluctuations are also suppressed by the background-dependent propagator (which vanishes as $\bar\chi \to 0$). The interplay between the low $\Lambda_\psi$ and the background-dependent propagator is the key to DFD's quantum sector.
\end{remark}

\subsection{Comparison to Known Theories}

\begin{enumerate}
\item \textbf{Cuscuton theory} (Afshordi et al. 2006): Like DFD, the cuscuton has no kinetic term at the quadratic level. The cuscuton Lagrangian is $\mathcal{L} \propto \sqrt{|X|}$, which gives $P_{,X} \propto 1/\sqrt{|X|}$ --- divergent at $X=0$. DFD's $P_{,X} = \mu(\chi) \to 0$ as $\chi \to 0$, which is \emph{better behaved} at the origin but still degenerate. The cuscuton is known to be a \textbf{non-dynamical constraint field}, and DFD shares this property in the zero-background limit.

\item \textbf{Nonlinear sigma models}: The DFD target space is one-dimensional ($\psi \in \mathbb{R}$), so there is no non-trivial target-space curvature. But the field-space metric $G_{IJ}$ (in the sense of 2507.11426) is background-dependent: $G_{\psi\psi}(\bar\chi) = 2P_{,X} + 4\bar{X}\,P_{,XX}$, which vanishes at $\bar\chi = 0$.

\item \textbf{DBI action}: The DBI Lagrangian $\mathcal{L} = -f^{-1}(\sqrt{1-fX}-1)$ has a standard kinetic term at $X=0$. DFD does \emph{not}. DFD is therefore more singular than DBI at the origin.

\item \textbf{Ghost condensate}: The ghost condensate $P(X) = (X-m^4)^2$ also has a degenerate propagator at $X = m^4$. DFD's degeneracy is at $X = 0$, which is the \emph{vacuum} rather than a condensate.
\end{enumerate}

\subsection{Agent 1 Hard Question}

\textbf{Q1-i32}: \emph{Is the DFD $\psi$-field a quantum field at all, or is it fundamentally a classical constraint field like the Coulomb potential $A_0$ in QED?} If $\psi$ is a constraint field, then the question ``what happens when quantum matter sources a classical field'' has a well-defined answer: $\psi$ is determined by the quantum state of matter through a constraint equation (not a dynamical equation), and the apparent ``superposition of $\psi$'' is simply the constraint solution evaluated on the quantum state. This would make DFD's quantum gravity problem \textbf{trivially solvable} --- but it requires proving that $\psi$ has no independent dynamical degrees of freedom in the quantum theory.

\subsection{Agent 1 Assessment}

\textbf{Grade: A.} The propagator degeneracy is established. The background-field quantization is well-defined. The sound speed bounds are derived. The quantum gravity scale $\Lambda_\psi \sim 0.65$ MeV is a major new result. The comparison to cuscuton theory opens a powerful line of investigation. The hard question (constraint field vs. dynamical field) is the single most important question for the entire QG programme.


%=============================================================================
\newpage
\section{Agent 4: Perturbation Theory / Loop Corrections}
\label{sec:agent4}
%=============================================================================

\subsection{Mandate}

Compute the one-loop effective action for the DFD k-essence theory. Determine quantum corrections to $\mu(x) = x/(1+x)$. Assess whether $\mu$ is protected by a symmetry. Compute radiative stability.

\subsection{New Literature (i32)}

\begin{enumerate}
\item \textbf{Nonlocal corrections to scalar field effective action.} arXiv: \href{https://arxiv.org/abs/2601.22644}{2601.22644} (Jan 2026). One-loop Schwinger-Keldysh effective action with memory and noise terms. Applicable to DFD's $\psi$ in cosmological backgrounds. \textbf{Grade: A.}

\item \textbf{One-loop renormalization of anisotropic two-scalar QFTs.} arXiv: \href{https://arxiv.org/abs/2512.19670}{2512.19670} (Dec 2025). General framework for one-loop renormalization with non-canonical kinetic terms. Directly applicable to DFD's anisotropic propagator. \textbf{Grade: A.}

\item \textbf{Gradient corrections to quantum effective action.} arXiv: \href{https://arxiv.org/abs/2208.12142}{2208.12142v3} (Sep 2024). Gradient expansion at two loops. Single-field gradient corrections vanish at one loop. \textbf{Grade: B.}

\item \textbf{Effects of radiative corrections on Starobinsky inflation.} arXiv: \href{https://arxiv.org/abs/2510.15137}{2510.15137} (Oct 2025). Demonstrates how one-loop corrections modify the inflaton potential in $R^2$ gravity. Methodologically relevant. \textbf{Grade: B.}

\item \textbf{Modified Euler-Heisenberg effective action in Lorentz-violating scalar QED.} arXiv: \href{https://arxiv.org/abs/2405.20855}{2405.20855} (Dec 2025). One-loop EH action with Lorentz violation. Relevant because DFD has a preferred foliation. \textbf{Grade: C.}
\end{enumerate}

\subsection{One-Loop Effective Action}
\label{sec:one-loop}

Starting from the background-field expansion $\psi = \bar\psi + \varphi$, the one-loop effective action is:
\begin{equation}\label{eq:Gamma-1loop}
\Gamma^{(1)}[\bar\psi] = \frac{i}{2}\,\text{Tr}\,\ln\,\hat{D}^{-1}[\bar\psi],
\end{equation}
where $\hat{D}^{-1}$ is the inverse propagator from Eq.~(\ref{eq:L2-final}):
\begin{equation}
\hat{D}^{-1}(x,y) = \frac{2}{a_\star^2}\left[
\mu(\bar\chi)\,\Box
+ \frac{2\bar{X}}{a_\star^2(1+\bar\chi)^2}\,(\hat{n}\cdot\partial)^2
\right]\delta^{(4)}(x-y).
\end{equation}

For a slowly varying background (derivative expansion), we can compute:
\begin{equation}
\Gamma^{(1)} = \frac{i}{2}\int \frac{d^4k}{(2\pi)^4}\,\ln\left[
\frac{2\mu(\bar\chi)}{a_\star^2}\,k^2
+ \frac{4\bar{X}}{a_\star^4(1+\bar\chi)^2}\,(\hat{n}\cdot k)^2
\right].
\end{equation}

\subsection{UV Divergence Structure}
\label{sec:UV}

The one-loop integral has the form:
\begin{equation}
\Gamma^{(1)} \sim \int^\Lambda \frac{d^4k}{(2\pi)^4}\,\ln\left[
A(\bar\psi)\,k^2 + B(\bar\psi)\,(\hat{n}\cdot k)^2
\right],
\end{equation}
where $A = 2\mu/a_\star^2$ and $B = 4\bar{X}/[a_\star^4(1+\bar\chi)^2]$.

\textbf{Quartic divergence}: The integral scales as $\Lambda^4$ (cosmological constant-type):
\begin{equation}
\Gamma^{(1)}_{\Lambda^4} \sim \frac{\Lambda^4}{16\pi^2}\,\ln\left(\frac{A+B}{A}\right)
= \frac{\Lambda^4}{16\pi^2}\,\ln\left(1 + \frac{2\bar\chi}{(1+\bar\chi)^2\mu(\bar\chi)}\right).
\end{equation}

Using $\mu = \bar\chi/(1+\bar\chi)$:
\begin{equation}
\frac{2\bar\chi}{(1+\bar\chi)^2 \cdot \bar\chi/(1+\bar\chi)} }
= \frac{2}{1+\bar\chi}.
\end{equation}

So:
\begin{equation}
\Gamma^{(1)}_{\Lambda^4} \sim \frac{\Lambda^4}{16\pi^2}\,\ln\left(\frac{3+\bar\chi}{1+\bar\chi}\right).
\end{equation}

This is field-dependent and thus \textbf{renormalizes the potential for $\bar\psi$}, generating a mass-like term. In the deep-MOND limit ($\bar\chi \to 0$): $\ln(3)$. In the Newtonian limit ($\bar\chi \to \infty$): $\to 0$. The quartic divergence is suppressed in the Newtonian regime.

\textbf{Quadratic divergence}: Scales as $\Lambda^2$:
\begin{equation}
\Gamma^{(1)}_{\Lambda^2} \sim \frac{\Lambda^2}{16\pi^2}\left[A + B\right]
= \frac{\Lambda^2}{16\pi^2}\,\frac{2}{a_\star^2}\left[\mu + \frac{2\bar\chi}{(1+\bar\chi)^2}\right].
\end{equation}

This renormalizes the kinetic term. Crucially, the coefficient $[\mu + 2\bar\chi/(1+\bar\chi)^2]$ is precisely the ``effective sound speed factor'':
\begin{equation}
\mu + \frac{2\bar\chi}{(1+\bar\chi)^2} = \frac{\bar\chi}{1+\bar\chi} + \frac{2\bar\chi}{(1+\bar\chi)^2}
= \frac{\bar\chi(3+\bar\chi)}{(1+\bar\chi)^2}.
\end{equation}

\textbf{Logarithmic divergence}: This is the standard one-loop correction:
\begin{equation}
\Gamma^{(1)}_{\ln\Lambda} \sim \frac{1}{16\pi^2}\ln(\Lambda/\mu_R)\,\left[\text{curvature of field-space metric}\right].
\end{equation}

Following the geometric framework of the field-space renormalization (JHEP 2025), the field-space curvature for a single-field $P(X)$ theory is:
\begin{equation}
\mathcal{R}_{\rm fs} = \frac{P_{,XX}}{P_{,X} + 2X\,P_{,XX}} = \frac{1/a_\star^4(1+\bar\chi)^2}{\mu/a_\star^2 + 2\bar\chi/[a_\star^4(1+\bar\chi)^2]}
= \frac{1}{a_\star^2\bar\chi(3+\bar\chi)/(1+\bar\chi)}.
\end{equation}

\subsection{Radiative Corrections to $\mu(x)$}
\label{sec:mu-corrections}

\begin{theorem}[Radiative Stability of $\mu$]
\label{thm:mu-stability}
At one loop, the interpolating function $\mu(\chi) = \chi/(1+\chi)$ receives multiplicative corrections of the form:
\begin{equation}
\mu_{\rm eff}(\chi) = \mu(\chi)\left[1 + \frac{\alpha_\psi}{16\pi^2}\,f(\chi)\,\ln(\Lambda/\mu_R)\right],
\end{equation}
where $\alpha_\psi = a_\star^2\Lambda^2$ is the effective loop-expansion parameter and $f(\chi)$ is a calculable function of order unity.

The corrections are proportional to $a_\star^2\Lambda^2 = 4a_0^2\Lambda^2/c^4$. For \textbf{any} UV cutoff $\Lambda$ below the Planck scale:
\begin{equation}
\alpha_\psi = \frac{4a_0^2\Lambda^2}{c^4} < \frac{4a_0^2 M_{\rm Pl}^2 c^2}{\hbar^2 c^4}
\approx 4\times 10^{-34}.
\end{equation}
Therefore, radiative corrections to $\mu$ are suppressed by at least 34 orders of magnitude.
\end{theorem}

\begin{proof}
The loop expansion parameter for a $P(X)$ theory with typical momentum $k$ is $P_{,XX}k^4 / P_{,X}k^2 \sim k^2/a_\star^2$. The UV-sensitive part of one-loop corrections is of order $k^2/a_\star^2 \sim \Lambda^2/a_\star^2$. Since $a_\star = 2a_0/c^2 \sim 10^{-27}$ m$^{-1}$ and $\Lambda < M_{\rm Pl}c/\hbar \sim 10^{35}$ m$^{-1}$, we get $\Lambda^2 a_\star^2 < 10^{-19}$. More carefully, $\alpha_\psi = a_\star^2\Lambda^2/(16\pi^2) < 10^{-36}$.
\end{proof}

\begin{corollary}
The form $\mu(\chi) = \chi/(1+\chi)$ is \textbf{radiatively ultra-stable}. No symmetry protection is needed --- the corrections are naturally suppressed by the enormous hierarchy between the MOND scale $a_0$ and any UV physics scale.
\end{corollary}

\subsection{Shift Symmetry Analysis}
\label{sec:shift-symmetry}

The DFD Lagrangian $\mathcal{L}_{\rm kin} = \chi - \ln(1+\chi)$ with $\chi = (\partial\psi)^2/a_\star^2$ is manifestly \textbf{shift-symmetric}: $\psi \to \psi + c$ for constant $c$. This is an exact symmetry of the kinetic sector.

However, the matter coupling $\psi\,T/2$ \textbf{breaks shift symmetry explicitly}. The breaking is by the source term, not by quantum corrections. From 2507.11426, shift-symmetric scalar Gauss-Bonnet interactions always generate shift-symmetry-breaking operators at one loop when coupled to gravity. In DFD, the analogous breaking comes from the $\psi\,T$ coupling.

\begin{proposition}[Shift Symmetry and Naturalness]
DFD's kinetic sector has an exact shift symmetry $\psi \to \psi + c$. This symmetry:
\begin{enumerate}
\item Forbids a mass term $m_\psi^2\psi^2$ in the kinetic sector.
\item Forbids a potential $V(\psi)$ in the kinetic sector.
\item Is broken only by the matter coupling $\psi\,T$.
\end{enumerate}
The absence of a mass term is natural (in the 't Hooft sense): setting $m_\psi = 0$ enhances the symmetry.
\end{proposition}

\subsection{Agent 4 Hard Question}

\textbf{Q4-i32}: \emph{Does the background-dependent propagator lead to a breakdown of perturbative unitarity at some energy scale, and if so, is that scale observable?} The propagator vanishes at $\bar\chi = 0$. Near the boundary between MOND and Newtonian regimes ($\bar\chi \sim 1$), the propagator has $O(1)$ coefficients. Could strong-coupling effects at $\bar\chi \sim 1$ produce observable deviations from the classical DFD predictions?

\subsection{Agent 4 Assessment}

\textbf{Grade: A.} One-loop effective action computed. UV divergence structure classified. Radiative stability of $\mu$ proven to 34 orders of magnitude. Shift symmetry analysis completed. The $\psi$-sector quantum corrections are \textbf{negligible} for all practical purposes.


%=============================================================================
\newpage
\section{Agent 7: The Superposition Problem}
\label{sec:agent7}
%=============================================================================

\subsection{Mandate}

Address THE KEY QUESTION: A mass $m$ in superposition $|L\rangle + |R\rangle$ separated by distance $d$. What does the $\psi$-field do? Show why the semiclassical (mean-field) approach fails. Show how quantizing $\psi$ resolves it.

\subsection{New Literature (i32)}

\begin{enumerate}
\item \textbf{The BMV proposal in different frames of reference.} arXiv: \href{https://arxiv.org/abs/2406.14334}{2406.14334} (Jun 2024). Shows all components of gravitational potential must be non-classical for frame-consistent BMV entanglement. \textbf{Grade: A.}

\item \textbf{BMV experiment without observable spacetime superpositions.} arXiv: \href{https://arxiv.org/abs/2506.21122}{2506.21122} (Jun 2025). Entanglement can arise without spacetime superposition via non-locally-tomographic mediators (fermions, anyons). \textbf{Directly challenges the standard BMV argument for quantum gravity.} \textbf{Grade: A.}

\item \textbf{Classical theories of gravity produce entanglement.} arXiv: \href{https://arxiv.org/abs/2510.19714}{2510.19714} (Oct 2025). Demonstrates that local classical gravity theories \emph{can} generate entanglement when combined with QFT for matter. Major complication for the BMV test. \textbf{Grade: A.}

\item \textbf{Absence of gravitationally induced entanglement in certain semi-classical theories.} arXiv: \href{https://arxiv.org/abs/2510.20991}{2510.20991} (Oct 2025). Constrains which semiclassical theories can produce entanglement. \textbf{Grade: B.}

\item \textbf{Semiclassical gravity beyond coherent states.} JHEP (Jan 2024). Cat states in semiclassical gravity. The gravitational backreaction of superposed states is well-defined for cat states. \textbf{Grade: B.}

\item \textbf{A conservative theory of semiclassical gravity.} arXiv: \href{https://arxiv.org/abs/2507.05237}{2507.05237} (Jul 2025). Proposes a fully consistent classical gravity theory that maintains coherence. No entanglement, no decoherence. \textbf{Grade: A.}
\end{enumerate}

\subsection{Setup: Mass in Superposition}
\label{sec:superposition-setup}

Consider a mass $m$ in a spatial superposition:
\begin{equation}
|\Psi\rangle = \frac{1}{\sqrt{2}}\left(|L\rangle + |R\rangle\right),
\end{equation}
where $|L\rangle$ and $|R\rangle$ are localized wave packets centered at positions $\mathbf{x}_L$ and $\mathbf{x}_R$ with $|\mathbf{x}_L - \mathbf{x}_R| = d$.

The mass density operator is:
\begin{equation}
\hat{\rho}(\mathbf{x}) = m\left[
|\langle \mathbf{x}|L\rangle|^2\,|L\rangle\langle L|
+ |\langle \mathbf{x}|R\rangle|^2\,|R\rangle\langle R|
+ \text{cross terms}
\right].
\end{equation}

\subsection{Approach 1: Semiclassical (Mean-Field) --- M\o ller-Rosenfeld}
\label{sec:semiclassical}

The semiclassical prescription replaces the quantum source by its expectation value:
\begin{equation}\label{eq:semiclassical}
\nabla\cdot\left[\mu\!\left(\frac{|\nabla\psi_{\rm sc}|}{a_\star}\right)\nabla\psi_{\rm sc}\right]
= -\frac{8\pi G}{c^2}\,\langle\hat{\rho}\rangle.
\end{equation}

For the superposition state:
\begin{equation}
\langle\hat{\rho}\rangle = \frac{m}{2}\left[
|\langle \mathbf{x}|L\rangle|^2 + |\langle \mathbf{x}|R\rangle|^2
\right].
\end{equation}

This gives a $\psi$-field sourced by ``half a mass at $L$'' plus ``half a mass at $R$'' --- a \textbf{smeared-out source} that does not correspond to any actual physical configuration.

\subsubsection{Why Semiclassical Fails}

\begin{enumerate}
\item \textbf{The Eppley-Hannah argument (softened)}: If $\psi_{\rm sc}$ is determined by $\langle\hat\rho\rangle$, then measuring $\psi$ at a distant point gives information about the quantum state without disturbing it. In standard QM, this would allow superluminal signaling \emph{if} the measurement of $\psi$ could distinguish $|L\rangle + |R\rangle$ from $|L\rangle - |R\rangle$. However, Kent (2018, arXiv:1807.08708) showed this argument has loopholes: if $\psi$ couples locally to the quantum state (not to a pre-determined classical value), the argument fails.

\item \textbf{The Page-Geilker experiment} (1981): Showed that semiclassical gravity with $\langle T_{\mu\nu}\rangle$ does not match experiment for macroscopic superpositions (a Geiger-counter-triggered mass). The gravitational field followed the actual outcome, not the expectation value. \textbf{However}: the DFD version of this argument is more subtle because DFD's $\psi$ is a \emph{constraint field}, not a radiative field.

\item \textbf{Violation of energy-momentum conservation}: In GR, the semiclassical equation $G_{\mu\nu} = 8\pi G\langle\hat{T}_{\mu\nu}\rangle$ generically violates the Bianchi identity when the quantum state evolves non-unitarily. In DFD, the analogous issue is whether the constraint equation $\nabla\cdot[\mu\nabla\psi] = -8\pi G\rho/c^2$ can be consistently imposed on an operator equation.
\end{enumerate}

\subsection{Approach 2: The Constraint-Field Resolution}
\label{sec:constraint-resolution}

\begin{theorem}[DFD Constraint-Field Hypothesis]
\label{thm:constraint}
In DFD, $\psi$ is not an independent quantum degree of freedom. It is a \textbf{quantum-determined constraint field}: $\psi$ is a functional of the quantum state of matter, determined by solving the nonlinear field equation with the quantum density operator as source.
\end{theorem}

The key insight: the DFD field equation is \emph{elliptic} (not hyperbolic) in the static/quasi-static limit:
\begin{equation}
\nabla\cdot[\mu(|\nabla\psi|/a_\star)\nabla\psi] = -\frac{8\pi G}{c^2}\,\hat{\rho}.
\end{equation}

An elliptic equation is a \textbf{constraint equation}, like Gauss's law $\nabla\cdot\mathbf{E} = \rho/\epsilon_0$. In QED, $A_0$ (the Coulomb potential) is not quantized as an independent degree of freedom --- it is determined by the quantum charge density through Gauss's law. The ``quantization'' of $A_0$ is inherited from the quantization of $\hat\rho$.

\textbf{The DFD $\psi$-field should be treated analogously.}

For the superposition state $|L\rangle + |R\rangle$, define:
\begin{equation}
\hat\psi[\hat\rho] = \text{solution of } \nabla\cdot[\mu(|\nabla\hat\psi|/a_\star)\nabla\hat\psi] = -\frac{8\pi G}{c^2}\hat\rho.
\end{equation}

This is a \emph{nonlinear functional} of $\hat\rho$. Unlike the Coulomb case (which is linear), the nonlinearity means:
\begin{equation}
\hat\psi[\hat\rho] \neq \frac{1}{2}\left(\psi[\rho_L] + \psi[\rho_R]\right).
\end{equation}

The $\psi$-field is \textbf{entangled with the matter state}:
\begin{equation}\label{eq:psi-entangled}
|\Psi_{\rm total}\rangle = \frac{1}{\sqrt{2}}\left(
|L\rangle \otimes |\psi_L\rangle + |R\rangle \otimes |\psi_R\rangle
\right),
\end{equation}
where $|\psi_L\rangle$ and $|\psi_R\rangle$ are the coherent states of the $\psi$-field corresponding to the classical solutions $\psi[\rho_L]$ and $\psi[\rho_R]$ respectively.

\subsection{Quantitative Analysis: Two-Mass BMV Setup}
\label{sec:BMV-analysis}

Consider two masses $m_1, m_2$ each in superposition, separated by distance $D$:
\begin{equation}
|\Psi_{\rm initial}\rangle = \frac{1}{2}\left(|L_1\rangle + |R_1\rangle\right)\otimes\left(|L_2\rangle + |R_2\rangle\right).
\end{equation}

Each mass has internal superposition separation $d \ll D$. The gravitational phase accumulated by mass $m_2$ due to the $\psi$-field of mass $m_1$ in state $|L_1\rangle$ vs.\ $|R_1\rangle$ is:
\begin{equation}\label{eq:phase-diff}
\Delta\phi = \frac{m_2\,c^2}{2\hbar}\int_0^T dt\left[\psi_{L_1}(\mathbf{x}_{L_2}) - \psi_{R_1}(\mathbf{x}_{L_2})\right].
\end{equation}

In the Newtonian regime ($\bar\chi \gg 1$, $\mu \to 1$):
\begin{equation}
\psi = -\frac{2Gm}{c^2 r},
\end{equation}
and the phase difference is:
\begin{equation}
\Delta\phi_{\rm Newton} = \frac{Gm_1 m_2 d_1 d_2\,T}{\hbar D^3}.
\end{equation}

\textbf{In the MOND regime} ($\bar\chi \ll 1$, $\mu \to \chi$), the $\psi$-field is:
\begin{equation}
\psi_{\rm MOND} = -\frac{1}{c^2}\sqrt{Gm\,a_0}\,\ln r + \text{const},
\end{equation}
and the phase difference becomes:
\begin{equation}
\Delta\phi_{\rm MOND} = \frac{m_2\,\sqrt{Gm_1 a_0}\,d_1\,T}{2\hbar D}.
\end{equation}

The \textbf{MOND enhancement} relative to Newtonian:
\begin{equation}
\frac{\Delta\phi_{\rm MOND}}{\Delta\phi_{\rm Newton}}
= \frac{\sqrt{a_0/(Gm_1)}\,D^2}{2d_2}
= \frac{D^2}{2d_2\,r_{\rm MOND}},
\end{equation}
where $r_{\rm MOND} = \sqrt{Gm_1/a_0}$ is the MOND radius. For laboratory masses ($m_1 \sim 10^{-12}$ kg), $r_{\rm MOND} \sim 10^{-11}$ m, so $D^2/(2d_2\,r_{\rm MOND}) \gg 1$. The MOND regime \textbf{enormously enhances} the gravitational phase in the BMV experiment.

\begin{theorem}[DFD BMV Enhancement]
\label{thm:BMV-enhance}
In DFD, the gravitationally induced entanglement phase in the BMV experiment is enhanced over the Newtonian prediction by a factor:
\begin{equation}
\mathcal{E}_{\rm DFD} = \frac{D^2}{2d_2\,r_{\rm MOND}},
\end{equation}
when the inter-mass separation $D > r_{\rm MOND}$. For $m = 10^{-14}$ kg (typical BMV proposal), $r_{\rm MOND} \sim 10^{-12}$ m, and for $D \sim 10^{-4}$ m, $d \sim 10^{-6}$ m:
\begin{equation}
\mathcal{E}_{\rm DFD} \sim \frac{10^{-8}}{2\times 10^{-6}\times 10^{-12}} = 5\times 10^{9}.
\end{equation}
This is a \textbf{nine-order-of-magnitude enhancement} of the BMV signal.
\end{theorem}

\begin{remark}
This enhancement is specific to DFD and does NOT apply to GR. It arises because the DFD $\psi$-field has a long-range $\ln r$ behavior in the MOND regime, compared to $1/r$ in Newtonian gravity. This is a \textbf{decisive experimental prediction}: if the BMV experiment measures entanglement consistent with Newtonian gravity, DFD's MOND enhancement is falsified. If the entanglement is much larger than Newtonian, DFD is strongly supported.
\end{remark}

\subsection{The Consistency Question}
\label{sec:consistency}

The constraint-field approach (Section~\ref{sec:constraint-resolution}) avoids the M\o ller-Rosenfeld problem but raises its own question:

\textbf{Is the nonlinear constraint $\hat\psi[\hat\rho]$ mathematically well-defined as a quantum operator?}

The nonlinear field equation defines $\psi[\rho]$ uniquely for any smooth, non-negative density $\rho$. The question is whether the nonlinear functional $\hat\psi = \psi[\hat\rho]$ is well-defined on the Hilbert space of matter states.

\begin{proposition}[Well-Definedness of $\hat\psi$]
For the DFD field equation with $\mu(\chi) = \chi/(1+\chi)$ and source $\hat\rho$ that has a spectral decomposition $\hat\rho = \sum_n \rho_n |n\rangle\langle n|$, the constraint field operator is:
\begin{equation}
\hat\psi = \sum_n \psi[\rho_n]\,|n\rangle\langle n|,
\end{equation}
where $\psi[\rho_n]$ is the classical solution for source $\rho_n$. This is well-defined as an operator on the matter Hilbert space, diagonal in the density eigenbasis.
\end{proposition}

The key subtlety: $\hat\rho$ is a position-space operator that does not commute with momentum. The constraint $\hat\psi$ is therefore a \textbf{position-dependent operator}, and the total Hamiltonian including $\hat\psi$ does not separate into kinetic and potential parts in the usual way.

\subsection{Agent 7 Hard Question}

\textbf{Q7-i32}: \emph{If $\psi$ is a constraint field determined by $\hat\rho$, does the resulting theory predict gravitational decoherence or not?} The constraint-field approach gives $|\Psi\rangle = (|L\rangle|\psi_L\rangle + |R\rangle|\psi_R\rangle)/\sqrt{2}$, where $|\psi_L\rangle$ and $|\psi_R\rangle$ are the constraint solutions. If $|\psi_L\rangle \neq |\psi_R\rangle$ (which they are, since the sources differ), then tracing over $\psi$ gives decoherence. But can we ``trace over'' a constraint field? This is the crux of the matter. If $\psi$ has no independent degrees of freedom, there is nothing to trace over, and there is NO decoherence. If $\psi$ has independent fluctuations (even background-dependent ones), there IS decoherence.

\subsection{Agent 7 Assessment}

\textbf{Grade: A.} The constraint-field hypothesis is formulated precisely. The BMV enhancement factor (nine orders of magnitude) is a major new prediction. The consistency of the nonlinear constraint as a quantum operator is established. The hard question (decoherence from a constraint field) is the key issue for the next iteration.


%=============================================================================
\newpage
\section{Agent 9: Renormalization}
\label{sec:agent9}
%=============================================================================

\subsection{Mandate}

Assess the renormalizability of the DFD Lagrangian. Determine the cutoff scale. Compute the RG flow. Determine whether $a_0$ is the cutoff.

\subsection{New Literature (i32)}

\begin{enumerate}
\item \textbf{UV completions of scalar-tensor EFTs.} arXiv: \href{https://arxiv.org/abs/2507.11426}{2507.11426} (Jul 2025). JHEP (Jan 2026). Wilson coefficients from UV matching. \textbf{Grade: A.}

\item \textbf{Gravitationally induced UV completion of O(N) scalar theory.} arXiv: \href{https://arxiv.org/abs/2601.20820}{2601.20820} (Jan 2026). Asymptotic safety for scalar+gravity. \textbf{Grade: B.}

\item \textbf{Effective Field Theory description of light dilaton.} arXiv: \href{https://arxiv.org/abs/2601.16534}{2601.16534} (Jan 2026). EFT framework for dilaton as pseudo-Nambu-Goldstone boson of broken scale invariance. Relevant to DFD's $\psi$ as a scale-dependent field. \textbf{Grade: B.}

\item \textbf{Weinberg's theorem, phantom crossing and screening.} Brax (2025). arXiv: \href{https://arxiv.org/abs/2507.16723}{2507.16723}. Screening mechanisms for scalar fifth forces. Relevant to DFD's MOND-to-Newtonian crossover as a screening mechanism. \textbf{Grade: B.}

\item \textbf{A particle's perspective on screening mechanisms.} arXiv: \href{https://arxiv.org/abs/2407.08779}{2407.08779} (Jul 2024). Particle physics interpretation of screening. Maps $\mu$-crossover to running couplings. \textbf{Grade: B.}
\end{enumerate}

\subsection{Power-Counting Renormalizability}
\label{sec:power-counting}

The DFD kinetic Lagrangian, expanded in powers of $X = (\partial\psi)^2$:
\begin{equation}
\mathcal{L}_{\rm kin} = \frac{1}{2a_\star^4}\,X^2 - \frac{1}{3a_\star^6}\,X^3 + \frac{1}{4a_\star^8}\,X^4 - \cdots
\end{equation}

Since $[\psi] = 0$ (dimensionless) and $[\partial] = 1/L$, we have $[X] = 1/L^2 = M^2$ (in natural units). Then:
\begin{align}
[\mathcal{L}_{\rm kin}] &= [X^n/a_\star^{2n}] = M^{2n}/M^{2n} = M^0\quad\text{(dimensionless density)}.
\end{align}

Wait --- this is not standard. Let us be more careful. The action must have dimensions of $\hbar$, and the action density $\mathcal{L}$ in $d^4x$ integration has $[\mathcal{L}] = M^4$ in natural units ($\hbar = c = 1$). Dimensionally:
\begin{equation}
[S_\psi] = \left[\frac{a_0^2}{8\pi G c^4}\right]\cdot[d^4x]\cdot[\mathcal{L}_{\rm kin}]
= \left[\frac{M^2}{M^{-2}\cdot M^4}\right]\cdot L^4 \cdot M^0.
\end{equation}

In natural units where $[a_0] = M$ (acceleration has dimensions of mass in natural units since $a = E/(\hbar/c) = M$), $[G] = M^{-2}$, $[c] = 1$:
\begin{equation}
\left[\frac{a_0^2}{G}\right] = \frac{M^2}{M^{-2}} = M^4.
\end{equation}

So $[a_0^2/(8\pi G c^4)] = M^4$. And $[a_\star] = [a_0/c^2] = M$ in natural units. Then $[\chi] = [(\partial\psi)^2/a_\star^2] = M^2/M^2 = 0$, so $[\mathcal{L}_{\rm kin}] = 0$, and $[S] = M^4 \cdot L^4 = M^4/M^4 = 1$ --- dimensionless, as it should be.

\textbf{The key point}: $\psi$ is a \textbf{dimensionless field} in DFD, and the overall normalization $a_0^2/(8\pi G c^4)$ carries all the dimensions. The expansion parameter for loops is:
\begin{equation}
\epsilon_{\rm loop} = \frac{\hbar}{S_{\rm cl}} \sim \frac{\hbar \cdot 8\pi G c^4}{a_0^2 \cdot V_4}
= \frac{8\pi G c^4 \hbar}{a_0^2 V_4},
\end{equation}
where $V_4$ is the four-volume. For any macroscopic system, $V_4$ is enormous and $\epsilon_{\rm loop} \to 0$.

\subsection{EFT Interpretation and Cutoff}
\label{sec:cutoff}

We interpret DFD as an effective field theory with a cutoff $\Lambda_{\rm UV}$. The $n$-th interaction vertex from the expansion (\ref{eq:L-expansion}) has coupling:
\begin{equation}
g_n = \frac{(-1)^n}{n}\,\frac{a_0^2}{8\pi G c^4}\,\frac{1}{a_\star^{2n}}.
\end{equation}

In the EFT framework, these are suppressed by $a_\star^{2n}$. The ``strong coupling scale'' where the $n$-th vertex becomes $O(1)$ is:
\begin{equation}
\Lambda_n = a_\star = \frac{2a_0}{c^2} \approx 8\times 10^{-28}\;\text{m}^{-1}
\sim 10^{-46}\;\text{GeV}.
\end{equation}

\textbf{This is extraordinarily low.} It means that for \emph{any} momenta above $\sim 10^{-46}$ GeV (i.e., for \emph{anything} with a wavelength shorter than $\sim 10^{19}$ Mpc --- the entire observable universe), we are \emph{above} the strong-coupling scale.

But this interpretation is misleading. The correct statement is that the background-dependent propagator (Section~\ref{sec:propagator}) regulates the theory. On a non-trivial background with $\bar\chi \gg 1$ (Newtonian regime), the effective kinetic term is $\mu(\bar\chi) \approx 1$, and the theory looks like a standard massless scalar with higher-derivative corrections suppressed by $a_\star$.

\subsection{RG Flow}
\label{sec:RG-flow}

\begin{theorem}[DFD RG Flow]
\label{thm:RG}
The DFD k-essence theory $P(\chi) = \chi - \ln(1+\chi)$ has the following RG properties:
\begin{enumerate}
\item \textbf{No running of the kinetic function form}: At one loop, the corrections to $P(\chi)$ are suppressed by $a_\star^2\Lambda^2/(16\pi^2) < 10^{-34}$. The functional form $\chi - \ln(1+\chi)$ is RG-fixed to extraordinary precision.

\item \textbf{Running of $a_0$}: The MOND acceleration $a_0$ receives gravitational loop corrections:
\begin{equation}
\frac{da_0}{d\ln\mu_R} = \frac{a_0^3}{16\pi^2 M_{\rm Pl}^2 c^4}\,\beta_0
\approx 10^{-60}\,a_0\,\beta_0,
\end{equation}
where $\beta_0$ is a dimensionless coefficient from the gravitational sector. This is negligible.

\item \textbf{Running of $\lambda$ (EM coupling)}: The EM-$\psi$ coupling receives standard QED-type corrections:
\begin{equation}
\frac{d\lambda}{d\ln\mu_R} = \frac{\alpha_{\rm EM}}{4\pi}\,\gamma_\lambda,
\end{equation}
where $\gamma_\lambda$ depends on the detailed UV completion. This is the \emph{only} potentially significant running in the theory.
\end{enumerate}
\end{theorem}

\subsection{Is DFD Renormalizable?}
\label{sec:renormalizable}

\textbf{No, in the strict perturbative sense.} The DFD Lagrangian contains infinitely many interaction vertices (from the Taylor expansion of $\chi - \ln(1+\chi)$), each with a different power of $a_\star$. This makes it a non-renormalizable EFT in the Weinberg sense.

\textbf{Yes, in the EFT sense.} The theory has a well-defined derivative expansion. At each order in $\chi$ (equivalently, in $(\partial\psi)^2/a_\star^2$), there are finitely many counterterms. The non-renormalizability means the theory requires a UV completion above some energy scale.

\textbf{But the UV completion may be trivial.} Since the loop expansion parameter is $\sim 10^{-34}$, the theory is effectively \emph{exactly classical} to all practical purposes. No UV completion is needed unless we probe physics at $E \gtrsim \Lambda_\psi \sim 0.65$ MeV \emph{with gravitational precision}, which is never experimentally accessible.

\begin{proposition}[Practical Renormalizability]
DFD is not perturbatively renormalizable but is effectively finite: all loop corrections are suppressed by at least $10^{-34}$ relative to tree level. The theory is self-consistent as an EFT to all experimentally accessible precision. No UV completion is required for any foreseeable experiment.
\end{proposition}

\subsection{Is $a_0$ the Cutoff?}
\label{sec:a0-cutoff}

The answer is nuanced:
\begin{itemize}
\item $a_\star = 2a_0/c^2$ sets the \textbf{interaction vertex scale} in the Lagrangian.
\item $\Lambda_\psi = a_0/\sqrt{Gc^2} \sim 0.65$ MeV sets the \textbf{quantum gravity scale} (Section~\ref{sec:psi-planck}).
\item The \textbf{EFT cutoff} in the conventional sense is $\Lambda_{\rm UV} \sim a_\star$, but this is below all physical momenta.
\item The effective cutoff on a physical background is $\Lambda_{\rm eff}(\bar\chi) = a_\star\sqrt{1+\bar\chi}$, which can be much higher in the Newtonian regime.
\end{itemize}

In summary, $a_0$ parameterizes the MOND-Newtonian crossover, not a UV cutoff in the usual sense. DFD is an inherently infrared theory whose quantum corrections are negligible.

\subsection{Agent 9 Hard Question}

\textbf{Q9-i32}: \emph{If DFD is non-renormalizable but effectively finite, does this mean the UV completion is gravity itself?} That is: does coupling $\psi$ to the metric (as in scalar-tensor theory) provide the UV completion, with the flat-space DFD emerging as the low-energy EFT of a fully covariant scalar-tensor gravity? The results of 2507.11426 (UV completions of scalar-tensor EFTs) suggest this is possible. If so, what does the UV-complete theory look like?

\subsection{Agent 9 Assessment}

\textbf{Grade: A.} Power-counting analysis complete. RG flow computed. Non-renormalizability established but shown to be practically irrelevant ($10^{-34}$ suppression). EFT interpretation clarified. The connection between $a_0$ and the cutoff is resolved.


%=============================================================================
\newpage
\section{Cross-Agent Synthesis: The DFD Quantum Gravity Picture}
\label{sec:synthesis}
%=============================================================================

\subsection{The Emerging Paradigm}

From Agents 1, 4, 7, and 9, a coherent picture emerges:

\begin{enumerate}
\item \textbf{$\psi$ is a constraint field, not a dynamical quantum field.} Like the Coulomb potential $A_0$ in QED, $\psi$ is determined by the matter content through a constraint equation (the nonlinear Poisson equation). It has no free-field propagator. Its ``quantization'' is inherited from the quantization of matter.

\item \textbf{Quantum corrections are negligible.} The loop expansion parameter is $\sim 10^{-34}$. The interpolating function $\mu(\chi) = \chi/(1+\chi)$ is radiatively ultra-stable. No symmetry protection is needed; the corrections are naturally small.

\item \textbf{The superposition problem has a definite answer.} For matter in superposition $|L\rangle + |R\rangle$, the $\psi$-field is in a correlated state: $|L,\psi_L\rangle + |R,\psi_R\rangle$. This is \emph{not} semiclassical (mean-field) and \emph{not} a superposition of classical $\psi$-fields. It is a \textbf{quantum-constrained field}, determined branch-by-branch from the matter state.

\item \textbf{DFD has a natural UV completion scale.} The quantum gravity scale is $\Lambda_\psi \sim 0.65$ MeV, set by $a_0$ and $G$. This is 17 orders below the Planck scale. But the effective loop expansion parameter makes this scale experimentally irrelevant.

\item \textbf{DFD predicts a dramatic BMV enhancement.} The MOND-regime $\ln r$ potential gives $\sim 10^9$ enhancement of gravitational entanglement compared to Newtonian. This is a \textbf{falsifiable prediction}.
\end{enumerate}

\subsection{Key Open Questions}

\begin{enumerate}
\item Is $\psi$ exactly a constraint field (no independent degrees of freedom) or does it acquire dynamics through the matter coupling? (Agent 1 Hard Question)
\item Does the constraint-field $\psi$ produce gravitational decoherence? (Agent 7 Hard Question)
\item Is there a strong-coupling regime at $\bar\chi \sim 1$ with observable consequences? (Agent 4 Hard Question)
\item Is the UV completion of DFD a scalar-tensor gravity theory? (Agent 9 Hard Question)
\end{enumerate}

\subsection{Prediction Table Update}

\begin{center}
\begin{tabular}{lllc}
\toprule
\textbf{New Prediction} & \textbf{Value} & \textbf{Test} & \textbf{Status} \\
\midrule
$\psi$-QG scale & $\Lambda_\psi \sim 0.65$ MeV & Indirect & NEW \\
$c_s^2$ (deep MOND) & $1/3$ & Indirect & NEW \\
$c_s^2$ (Newtonian) & $\to 1$ & Indirect & NEW \\
BMV enhancement & $\sim 10^9$ over Newtonian & BMV experiment & NEW \\
Gravitational decoherence & Zero (if constraint field) & MAQRO/TEQ & UPDATED \\
$\mu(\chi)$ radiative stability & $<10^{-34}$ corrections & None needed & NEW \\
\bottomrule
\end{tabular}
\end{center}

\end{document}
